\section{Componente di generazione (LLM)}\label{sec:generazione}

Il componente Generation riceve la domanda e i frammenti selezionati dal recupero e produce la risposta mostrata allo
studente, con l'elenco delle fonti da cui è tratta. È l'unico punto del sistema in cui il testo viene scritto anziché
cercato, e quindi l'unico in cui un'informazione può essere inventata: le scelte descritte in questa sezione ---
il modello, le istruzioni che riceve, il comportamento quando non ha elementi per rispondere --- sono orientate a ridurre questo rischio e a renderlo verificabile, coerentemente con il vincolo di
precisione fattuale introdotto all'inizio del capitolo.

\subsection{Scelta del modello}\label{sec:scelta-modello}

Il modello di produzione è Qwen2.5-14B-Instruct, scelto tra i quattro candidati presentati nella
Sezione~\ref{sec:llm} (Tabella~\ref{tab:llm_comparison}), tutti in versione quantizzata a 4 bit (AWQ) per
rientrare nei 16~GB della GPU dedicata e serviti da vLLM~\cite{kwon2023vllm}. La scelta si basa su un confronto
sperimentale tra i candidati, in cui è il modello che risponde in modo più corretto;
il prezzo è una latenza maggiore rispetto ai candidati più piccoli, ritenuta accettabile a fronte del vincolo di
precisione fattuale. Poiché il modello è servito da vLLM con un nome logico fisso, sostituirlo con uno dei candidati
alternativi richiede solo di cambiare la configurazione.

Il modello genera a temperatura zero, cioè scegliendo sempre la parola più probabile: in un compito in cui la
risposta deve riportare fedelmente un contenuto dato, la variabilità introdotta dal campionamento non porta alcun
beneficio, e la generazione deterministica rende riproducibili sia le risposte sia le valutazioni.

\subsection{Struttura del prompt}\label{sec:prompt}

Il modello riceve una conversazione composta da tre parti (Figura~\ref{fig:prompt}): un messaggio di sistema con le
istruzioni, fisso per ogni domanda; i turni precedenti della conversazione, fino agli ultimi cinque; e il messaggio
dello studente, che il sistema costruisce unendo i frammenti recuperati, il corso dello studente, se noto, e la
domanda.

\begin{figure}[H]
  \centering
  \begin{tikzpicture}[
      part/.style={draw, thick, rounded corners=3pt, align=left, text width=13.2cm, font=\scriptsize, inner sep=5pt},
      tag/.style={font=\scriptsize\bfseries, anchor=south west, inner sep=1pt},
      note/.style={font=\scriptsize\itshape, text=black!65, anchor=west, align=left, text width=3.0cm},
    ]
    \node[part, fill=gray!10] (sys) at (0,0) {%
      Sei un assistente virtuale per Università di Bologna\dots{} Rispondi sempre in italiano. [\dots]
      Per le domande pertinenti, usa esclusivamente le informazioni presenti nel CONTESTO fornito. Se il contesto
      non contiene la risposta, dillo esplicitamente invece di inventare\dots{} Riporta date, importi e numeri
      esattamente come appaiono nel contesto\dots{} [\dots] Per ogni affermazione cita il numero del blocco da cui
      l'hai presa\dots{} Non scrivere URL, link o nomi di file\dots};
    \node[tag] at (sys.north west) {Messaggio di sistema (istruzioni)};

    \node[part, fill=yellow!8, below=7mm of sys] (hist) {%
      \textbf{Studente:} \dots\quad \textbf{Assistente:} \dots\qquad (ultimi cinque turni, senza i numeri di
      citazione)};
    \node[tag] at (hist.north west) {Cronologia della conversazione};

    \node[part, fill=blue!6, below=7mm of hist] (user) {%
      \textbf{CONTESTO:}\\[2pt]
      \textbf{[1]} (corso: Infermieristica) Frequenza delle lezioni\\
      La frequenza ai corsi è obbligatoria per il 66\% delle ore di lezione. La verifica della frequenza viene
      effettuata con varie modalità definite dall'Ateneo\dots\\[1pt]
      ---\\[1pt]
      \textbf{[2]} (generale di ateneo) Lezioni: frequenza e orari\\
      Le modalità in cui si svolgono le lezioni in un Corso di Laurea\dots{} Ogni corso può prevedere la frequenza
      obbligatoria delle lezioni o meno\dots\\[1pt]
      --- \quad \dots{} \quad (altri blocchi, dal più pertinente al meno pertinente)\\[4pt]
      \textbf{CORSO DELLO STUDENTE:} Infermieristica\\[4pt]
      \textbf{DOMANDA (STUDENTE):}\\
      quante lezioni devo seguire per forza?};
    \node[tag] at (user.north west) {Messaggio dello studente, costruito dal sistema};
  \end{tikzpicture}
  \caption{Struttura del prompt inviato al modello, con un esempio reale: la domanda ``quante lezioni devo seguire
  per forza?'' posta da uno studente di Infermieristica. Ogni frammento è un blocco numerato, preceduto
  dall'indicazione del corso o dell'ateneo e dal titolo della pagina e della sezione da cui proviene.}
  \label{fig:prompt}
\end{figure}

\paragraph{I frammenti.} Ogni frammento diventa un blocco numerato, intestato con il titolo della pagina e della
sezione da cui proviene e con il suo ambito, il corso di laurea o l'indicazione ``generale di ateneo''. Il numero
permette al modello di indicare il blocco da cui trae ogni affermazione, e al sistema di trasformare quei numeri
nell'elenco delle fonti mostrato sotto la risposta, con i titoli e gli indirizzi reali delle pagine; l'ambito permette al modello di distinguere informazioni che,
sulle pagine di corsi diversi, hanno lo stesso titolo ma contenuto diverso. I blocchi sono disposti dal più
pertinente al meno pertinente e la domanda è posta in fondo, dopo di essi: in questo modo il frammento più
pertinente e la domanda occupano le due posizioni del contesto che, come discusso nella
Sezione~\ref{sec:rag} (Figura~\ref{fig:u-curve}), il modello sfrutta meglio. Con pochi blocchi il contesto
resta comunque breve, e l'effetto di un'informazione persa nel mezzo è contenuto.

\paragraph{La domanda.} Il modello non risponde al testo digitato dallo studente, ma alla sua versione condensata
(Sezione~\ref{sec:analisi-domanda}), la stessa usata per il recupero: in questo modo non deve ricostruire da solo a
cosa si riferisca una domanda mal formata rispetto a frammenti recuperati per la domanda esplicita. Quando la
domanda è stata scomposta in più parti, il modello le riceve anche separate e numerate, con l'istruzione di
rispondere a ciascuna: la sola versione condensata può infatti confonderle.

\paragraph{Le istruzioni.} Il messaggio di sistema traduce in regole il vincolo di precisione fattuale e i problemi
osservati durante lo sviluppo. Le principali sono:
\begin{itemize}
  \item \textbf{usare solo il contesto}: le informazioni vanno tratte esclusivamente dai blocchi forniti e, se la
  risposta non c'è, il modello deve dirlo invece di ricorrere alla propria conoscenza generale, che per scadenze e
  importi di un ateneo specifico è inaffidabile per costruzione (Sezione~\ref{sec:problema});
  \item \textbf{riportare i numeri alla lettera}: date, importi e soglie vanno trascritti esattamente come
  compaiono, senza arrotondamenti o parafrasi;
  \item \textbf{non riusare le risposte precedenti}: un'informazione data in un turno precedente vale per la domanda
  attuale solo se è presente anche nel contesto di questo turno, perché in una conversazione che cambia argomento o
  corso il modello tendeva a riproporre dati riferiti a un caso diverso;
  \item \textbf{elencare tutte le casistiche}: se il contesto riporta risposte diverse per situazioni diverse ---
  sessioni di laurea, rate, anni di corso --- e la domanda non specifica la propria, il modello deve elencarle tutte,
  ciascuna con la situazione a cui si riferisce, invece di sceglierne una arbitrariamente;
  \item \textbf{rispettare il corso dello studente}: se il corso è indicato, vanno usati solo i blocchi di quel corso
  e quelli di ateneo; se l'informazione esiste solo per un altro corso, il modello può riportarla, ma dichiarando a
  quale corso si riferisce. Senza questa indicazione, a una domanda sulla prova finale di Ingegneria e scienze
  informatiche il modello aveva risposto con le scadenze di Psicologia, presenti nel contesto con un titolo
  identico;
  \item \textbf{rivolgersi direttamente allo studente}: la risposta non deve nominare ``il contesto'' o ``i blocchi'',
  che sono dettagli del funzionamento interno, e se un'informazione manca deve dire semplicemente che non è
  disponibile nei documenti consultati;
  \item \textbf{citare le fonti}: ogni affermazione va accompagnata dal numero del blocco da cui è tratta.
\end{itemize}

\subsection{Gestione dei casi particolari}\label{sec:casi-particolari}

Non tutti i messaggi richiedono, o permettono, una risposta generata dai documenti. La
Figura~\ref{fig:casi-particolari} riassume come il sistema li distingue.

\begin{figure}[H]
  \centering
  \begin{tikzpicture}[
      dec/.style={draw, thick, rounded corners=2pt, fill=yellow!10, align=center, font=\scriptsize,
                  text width=2.6cm, minimum height=0.9cm},
      act/.style={draw, thick, rounded corners=3pt, fill=blue!6, align=center, font=\scriptsize,
                  text width=3.0cm, minimum height=0.9cm},
      fix/.style={act, fill=gray!10},
      arr/.style={-{Stealth[length=2mm]}, thick, black!70},
      yes/.style={font=\scriptsize\itshape, text=black!60},
    ]
    \node[act] (msg) at (0,0) {Messaggio dello studente};
    \node[act] (an) at (0,-1.5) {Condensazione e scomposizione (Sezione~\ref{sec:analisi-domanda})};
    \node[dec] (kind) at (0,-3.2) {È una domanda sull'università?};
    \node[fix] (smallr) at (5.0,-2.4) {Risposta di cortesia scritta dal modello, nessuna ricerca};
    \node[fix] (topicr) at (5.0,-4.0) {Messaggio fisso: l'assistente risponde solo su temi universitari};
    \node[act] (ret) at (0,-4.9) {Recupero dei frammenti (Sezione~\ref{sec:retrieval})};
    \node[dec] (found) at (0,-6.5) {Frammenti pertinenti?};
    \node[act] (gen) at (5.0,-6.5) {Risposta generata dai frammenti, con le fonti};
    \node[act] (course) at (0,-8.3) {Richiesta del corso o, se già tentata, messaggio fisso di informazione
      non trovata};

    \draw[arr] (msg) -- (an);
    \draw[arr] (an) -- (kind);
    \draw[arr] (kind.east) -- node[yes, above, sloped] {saluto} (smallr.west);
    \draw[arr] (kind.east) -- node[yes, below, sloped] {estranea} (topicr.west);
    \draw[arr] (kind) -- node[yes, right] {sì} (ret);
    \draw[arr] (ret) -- (found);
    \draw[arr] (found) -- node[yes, above] {sì} (gen);
    \draw[arr] (found) -- node[yes, right] {no} (course);
  \end{tikzpicture}
  \caption{Come il sistema tratta un messaggio. I messaggi che non sono domande sull'università vengono riconosciuti
  prima del recupero, dalla stessa chiamata al modello che scompone la domanda; il modello genera una risposta dai
  documenti solo quando il recupero ha trovato frammenti pertinenti.}
  \label{fig:casi-particolari}
\end{figure}

\paragraph{Messaggi che non sono domande.} I saluti e i ringraziamenti (``ciao'', ``grazie mille'') e le domande
estranee all'università (``chi ha vinto la Serie A?'') non devono arrivare al recupero. Cercati come domande, quasi
sempre non trovano nulla, e il sistema chiederebbe allo studente il suo corso di laurea; talvolta trovano invece un
frammento sopra soglia (``buongiorno'' recupera una pagina sulle scuole di specializzazione), che comparirebbe tra
le fonti della risposta.

Il riconoscimento avviene nella scomposizione (Sezione~\ref{sec:analisi-domanda}), senza chiamate aggiuntive al
modello. Prima di scomporre la domanda, il modello stabilisce se il messaggio è un saluto, a cui risponde con una
breve frase cordiale, o se è estraneo all'università, caso in cui il sistema risponde con un messaggio fisso; in
entrambi i casi non si esegue alcuna ricerca. Nel dubbio il messaggio è trattato come una domanda, e lo stesso vale
per un saluto seguito da una domanda (``grazie, e la seconda rata?''). La condensazione, infine, lascia invariati i
messaggi di cortesia, che altrimenti verrebbero riformulati come la domanda precedente.

\paragraph{Nessun frammento pertinente.} Quando il recupero non trova frammenti sopra soglia, nemmeno dopo i
tentativi descritti nella Sezione~\ref{sec:retrieval}, il modello non viene interpellato: la risposta è un messaggio
fisso, che dichiara di non aver trovato l'informazione e suggerisce di rivolgersi al sito o all'ufficio competente.
La scelta di non lasciare la decisione al modello è dettata dall'esperienza: senza contesto, istruito a dire che
l'informazione manca, il modello preferiva comunque formulare una risposta plausibile, arrivando a inventare la
citazione di una fonte inesistente. Prima di arrivare a questo messaggio il sistema chiede allo studente il corso
di laurea, quando la risposta potrebbe trovarsi tra le pagine di un corso (Sezione~\ref{sec:corso-priorita}).

\paragraph{Informazione parzialmente presente.} Il caso più delicato è quello in cui il recupero trova frammenti
pertinenti all'argomento, ma nessuno contiene l'informazione richiesta: una domanda sugli orari di apertura della
biblioteca del campus recupera pagine sulle biblioteche e sui servizi, nessuna delle quali riporta gli orari. Qui la
decisione spetta al modello, che in base alle istruzioni deve dichiarare che l'informazione non è disponibile.
Nelle prove il modello adottato si comporta in questo modo, rispondendo ad esempio che gli orari ``non sono
presenti'' e suggerendo di contattare la biblioteca; talvolta però, contro le istruzioni, si riferisce al
``contesto fornito'', esponendo allo studente un dettaglio del funzionamento interno.
